% ============================================================================
% Fichier  : partie3/P03_C10_admissibilite_preuve.tex
% Titre    : L'Admissibilité de la Preuve Numérique
% Auteur   : MINKA MI NGUIDJOÏ Thierry Emmanuel
% ============================================================================
% Sources : ISO/IEC 27037 ; RFC 3227 ; loi 2010/012. Opération FARO : fictif.
% ============================================================================

\chapter{L'Admissibilité de la Preuve Numérique}
\label{chap:preuve}
\chaptermark{L'admissibilité de la preuve}

\epigraph{%
    « Une preuve que l'on ne peut pas défendre\\
    ne vaut pas mieux qu'une preuve que l'on n'a pas. »%
}{Principe du cours}

\niveauChapitre{Intermédiaire}{Chapitres~\ref{chap:normes} et~\ref{chap:droit}.}

\begin{boxobjectifs}
\begin{itemize}
    \item Définir l'admissibilité et ses conditions.
    \item Comprendre l'intégrité, l'authenticité et la traçabilité.
    \item Relier l'admissibilité au Trilemme CRO.
    \item Reconnaître ce qui détruit une preuve.
    \item Préparer une preuve opposable dès l'acquisition.
\end{itemize}
\end{boxobjectifs}

\bigskip

% ============================================================================
\section*{Question directrice : la preuve recueillie chez LiZbiZ tiendrait-elle devant un tiers~?}
% ============================================================================

Recueillir une trace est une chose~; pouvoir s'en prévaloir devant la
direction, l'ANTIC ou un juge en est une autre. Ce chapitre pose les
conditions de l'admissibilité et montre comment une manipulation imprudente
détruit une preuve pourtant exacte.

\medskip
\textbf{Plan du chapitre.}
\begin{enumerate}
    \item Qu'est-ce que l'admissibilité (section~\ref{sec:c10-definir}).
    \item Intégrité, authenticité, traçabilité (section~\ref{sec:c10-conditions}).
    \item Le Trilemme CRO appliqué (section~\ref{sec:c10-cro}).
    \item Ce qui détruit une preuve (section~\ref{sec:c10-detruire}).
\end{enumerate}

\separateur

% ============================================================================
\section{Qu'est-ce que l'Admissibilité}
\label{sec:c10-definir}
% ============================================================================

\begin{boxdef}
\termecle{Admissibilité}~: aptitude d'une preuve à être reçue et prise en
compte par un tiers (direction, autorité, juge). Elle ne tient pas à la seule
exactitude de la trace, mais à la \emph{manière} dont elle a été recueillie,
conservée et documentée.
\end{boxdef}

% ============================================================================
\section{Intégrité, Authenticité, Traçabilité}
\label{sec:c10-conditions}
% ============================================================================

\begin{tableaucours}{L{3.2cm}L{9.0cm}}{%
    \tableauHeader{Condition} &
    \tableauHeader{Ce qu'elle exige}
}
Intégrité & La preuve n'a pas été altérée depuis l'acquisition (condensat). \\
Authenticité & La preuve est bien ce qu'elle prétend être, de la source déclarée. \\
Traçabilité & Chaque manipulation est documentée (chaîne de custody). \\
Licéité & La collecte respecte le droit (autorisation, proportionnalité). \\
\end{tableaucours}
\vspace{-0.8em}
\begin{center}{\small\itshape Tableau~10.1. Conditions de l'admissibilité}\end{center}

% ============================================================================
\section{Le Trilemme CRO Appliqué}
\label{sec:c10-cro}
% ============================================================================

\begin{boxCRO}
\textbf{L'admissibilité est l'opposabilité.} La dimension O du Trilemme est
précisément l'admissibilité~: une preuve opposable est une preuve admissible.
La F (fiabilité) recouvre l'intégrité et l'authenticité~; la C
(confidentialité) impose de ne saisir que le nécessaire. Une bonne pratique
vise \cro{0}{+}{+} à chaque acte.
\end{boxCRO}

% ============================================================================
\section{Ce qui Détruit une Preuve}
\label{sec:c10-detruire}
% ============================================================================

\begin{boxalert}
\textbf{Violation typique.} Travailler directement sur l'original, sans
condensat ni copie.

\smallskip
\textbf{Règle~:} on travaille sur une copie vérifiée par condensat, on
documente chaque acte, on ne modifie jamais l'original. Une preuve manipulée
sans trace est perdue, même si son contenu est exact.
\end{boxalert}

\begin{boxscenario}{Opération FARO. La preuve sauvée de justesse}
\textbf{Constatation.} Un technicien de LiZbiZ, de bonne foi, avait
redémarré le serveur \texttt{FS-01} et commencé à « nettoyer » avant
l'arrivée de M.~BIKORO. La mémoire vive était perdue~; le disque, lui, n'avait
pas encore été écrasé.

\medskip
\textbf{Analyse} (confiance élevée). Le redémarrage a détruit une preuve
volatile (mémoire) et a failli compromettre le disque. M.~BIKORO a figé
l'état restant, calculé les condensats et documenté la chaîne~: le disque
reste opposable, la mémoire est perdue. La leçon~: ne rien toucher avant
d'avoir figé.
\end{boxscenario}

\begin{boxmetier}{Analyste forensique}
\textbf{Figer avant de comprendre.} L'analyste résiste à la tentation
d'« aller voir » sur le système vivant~: chaque action non documentée dégrade
l'opposabilité. On fige, on copie, on vérifie, puis on analyse la copie.
\end{boxmetier}

\separateur

\begin{boxresume}
\textbf{Ce qu'il faut retenir de ce chapitre~:}
\begin{itemize}
    \item \textbf{L'admissibilité} tient à la manière, pas à la seule
          exactitude.
    \item \textbf{Intégrité, authenticité, traçabilité, licéité}~: les
          conditions.
    \item \textbf{L'opposabilité} du CRO \emph{est} l'admissibilité.
    \item \textbf{Travailler sur l'original} détruit la preuve.
    \item \textbf{Figer avant de comprendre.}
\end{itemize}
\end{boxresume}

\begin{boxinfo}
\textbf{Points de contrôle de l'élève.} Je sais~: définir l'admissibilité~;
citer ses conditions~; la relier au CRO~; reconnaître ce qui détruit une
preuve~; préparer une preuve opposable.
\end{boxinfo}

\bigskip

\begin{boxexo}[title={\ding{45}\quad Exercice~10.1. Définir \niveaubase}]
L'admissibilité tient-elle à l'exactitude de la trace ou à la manière de la
recueillir~?
\end{boxexo}

\begin{boxexo}[title={\ding{45}\quad Exercice~10.2. Conditions \niveaubase}]
Citez les quatre conditions de l'admissibilité.
\end{boxexo}

\begin{boxexo}[title={\ding{45}\quad Exercice~10.3. Condensat \niveaubase}]
À quoi sert le condensat (empreinte) d'une preuve~?
\end{boxexo}

\begin{boxexo}[title={\ding{45}\quad Exercice~10.4. CRO \niveaumoyen}]
Pourquoi dit-on que l'opposabilité du CRO est l'admissibilité~?
\end{boxexo}

\begin{boxexo}[title={\ding{45}\quad Exercice~10.5. FARO \niveaumoyen}]
Quelle preuve le redémarrage a-t-il détruite~? Laquelle a été sauvée~?
\end{boxexo}

\begin{boxexo}[title={\ding{45}\quad Exercice~10.6. L'original \niveaumoyen}]
Pourquoi ne jamais travailler sur l'original~?
\end{boxexo}

\begin{boxexo}[title={\ding{45}\quad Exercice~10.7. Bonne foi \niveaumoyen}]
Le technicien était de bonne foi. Pourquoi son geste était-il néanmoins
dommageable~?
\end{boxexo}

\begin{boxexo}[title={\ding{45}\quad Exercice~10.8. Licéité \niveaumoyen}]
En quoi une collecte illicite compromet-elle l'admissibilité~?
\end{boxexo}

\begin{boxexo}[title={\ding{45}\quad Exercice~10.9. Note de première intervention \niveauexpert}]
Rédigez la consigne de première intervention à afficher près des serveurs de
LiZbiZ (« que faire, que ne pas faire »).
\end{boxexo}

\begin{boxexo}[title={\ding{45}\quad Exercice~10.10. Dossier de preuve \niveauexpert}]
Décrivez le dossier qui accompagne une preuve pour la rendre opposable.
\end{boxexo}

\bigskip

\begin{boxinfo}
\textbf{Pour aller plus loin.}
\begin{itemize}
    \item ISO/IEC 27037~\cite{iso27037}~; RFC~3227~\cite{rfc3227}.
\end{itemize}
\end{boxinfo}

\begin{boxinfo}
\textbf{Transition.} La Partie~IV entre dans la pratique de l'enquête~: le
Chapitre~\ref{chap:fondements} pose les fondements du forensique et la chaîne
de custody.
\end{boxinfo}
